\section{Products}\label{sec:mp-products}
Each product with its payoff and how it is priced: in closed form on the market data of \S\ref{sec:mp-market}, or by
Monte Carlo on a model of \S\ref{sec:mp-models}. Sensitivities are derivatives of the price w.r.t.\ the market quotes:
delta and gamma on the spot, vega per vol node, rho per curve pillar, sensitivities to the dividend amounts; theta is
the derivative w.r.t.\ the time to maturity with rates and dividend inputs held.

% ─────────────────────────────────────────────────────────────────────────────
\subsection{Linear rate products}\label{sec:mp-linear}
Notation: $P_d$ the discount (overnight) curve of the currency, $P_f$ the projection curve of the index, $N$ the
notional, $\tau$ the accrual fraction.
\begin{itemize}[nosep]
  \item \emph{Deposit}: $\mathrm{NPV} = N\,P_d(0,s) - N(1 + r\tau)\,P_d(0,e)$.
  \item \emph{FRA}: $\mathrm{NPV} = \pm N\tau\,(F - K)\,P_d(0,T_1)/(1 + F\tau)$, $F$ the forward of the index on $P_f$, settled at the
        start; zero once started.
  \item \emph{Fixed leg}: $\sum_i N K \tau_i\,P_d(0,p_i)$.
  \item \emph{Floating leg} (Ibor): coupon $N(g\,F_i + m)\,\tau_i$ discounted at $P_d(0,p_i)$, gearing $g$, spread $m$,
        forward $F_i = \big(P_f(0,s_i)/P_f(0,e_i) - 1\big)/\tau^{\mathrm{idx}}_i$.
  \item \emph{Swap}: payer $=$ floating $-$ fixed leg; par rate $=$ floating leg$/$annuity.
  \item \emph{OIS}: compounded overnight coupon $\big(P_d(0,s_i)/P_d(0,e_i) - 1\big)/\tau_i$ on the discount curve itself, so
        the floating leg telescopes to $N\big(P_d(0,s) - P_d(0,e)\big)$.
\end{itemize}

% ─────────────────────────────────────────────────────────────────────────────
\subsection{Caps, floors and collars}\label{sec:mp-capfloor}
A cap (floor) is a strip of calls (puts) on the Ibor fixings, the first period excluded:
\[
  \mathrm{caplet}_i = N\tau_i\,P_d(0,p_i)\,\mathrm{Opt}(F_i, K, \sigma_i, T^{\mathrm{fix}}_i),
\]
$\mathrm{Opt}$ the formula of \citet{bachelier1900} on normal caplet vols, of \citet{black1976} on lognormal ones, $\sigma_i$ the caplet
vol at $(T^{\mathrm{fix}}_i, K)$; a caplet already fixed is the intrinsic value on the fixing, a paid one is dropped.
Collar $=$ cap $-$ floor. Under a Gaussian short-rate model the caplets have their own closed form (\S\ref{sec:mp-gsr}),
used in its calibration.

% ─────────────────────────────────────────────────────────────────────────────
\subsection{Swaptions}\label{sec:mp-swaption}
European, physically settled, into a fixed-vs-Ibor swap. With the annuity of the fixed leg $A = \sum_i N\tau_i\,P_d(0,p_i)$ and
the forward swap rate $F = \mathrm{PV}_{\mathrm{float}}/A$,
\[
  V = A\,\mathrm{Opt}(F, K, \sigma, T),\qquad \sigma = \sigma_{\mathrm{cube}}(T, \text{tenor}, K - F),
\]
Bachelier (normal cube) or Black-76 (lognormal); see \citet{brigo2006}. Model prices under G1++ / G2++: \S\ref{sec:mp-gsr}.

% ─────────────────────────────────────────────────────────────────────────────
\subsection{Vanilla options}\label{sec:mp-vanilla}
Calls and puts on an equity, $\phi = \pm 1$, payoff $(\phi(S_T - K))^+$; European or American.

\paragraph{Black--Scholes} \citep{black1973,merton1973}. European on the forward $F$ of \S\ref{sec:mp-forward} and the
surface vol at $(K, T)$:
\[
  V = P(0,T)\,\phi\,\big(F N(\phi d_1) - K N(\phi d_2)\big),\qquad d_{1,2} = \frac{\ln(F/K) \pm \tfrac12 \sigma^2 T}{\sigma\sqrt T}.
\]
Two dividend models \citep{haug2003}: \emph{discrete}, on the escrowed spot $S_0 - \mathrm{PV}_{\mathrm{div}}(T)$ with carry $b = r - q_r$;
\emph{continuous}, with the dividends replaced by the equivalent yield $q_{\mathrm{eff}}(T)$, held fixed when the spot moves
(sticky yield), carry $b = r - q_{\mathrm{eff}} - q_r$. The two give the same price and different deltas. With only a yield in
the data, $b = r - q - q_r$.

\paragraph{American} \citep{bjerksund1993}, see also \citet{haug2007}. A call with carry $b$ is exercised at a flat
boundary $I$:
\begin{gather*}
  I = B_0 + (B_\infty - B_0)\big(1 - e^{h(T)}\big),\qquad h(T) = -\big(bT + 2\sigma\sqrt T\big)\frac{B_0}{B_\infty - B_0},\\
  B_\infty = \frac{\beta}{\beta - 1}K,\qquad B_0 = \max\Big(K, \frac{r}{r-b}K\Big),\qquad
  \beta = \tfrac12 - \frac{b}{\sigma^2} + \sqrt{\Big(\frac{b}{\sigma^2} - \tfrac12\Big)^2 + \frac{2r}{\sigma^2}},
\end{gather*}
value $\alpha S^\beta - \alpha\,\phi(S,T,\beta,I,I) + \phi(S,T,1,I,I) - \phi(S,T,1,K,I) - K\phi(S,T,0,I,I) + K\phi(S,T,0,K,I)$ with
$\alpha = (I - K)I^{-\beta}$; $S - K$ above $I$; the European value when $b \ge r$. Puts by the transformation
$P(S, K, r, b) = C(K, S, r - b, -b)$.

\paragraph{Monte Carlo} (European). Paths of one of the models of \S\ref{sec:mp-models}, calibrated to the market, on a
uniform time grid ending at the expiry; price $P(0,T)$ times the mean payoff, with its standard error. The normals are
drawn from a fixed seed, so a repricing with bumped market data uses the same random numbers.

% ─────────────────────────────────────────────────────────────────────────────
\subsection{Digital options}\label{sec:mp-digital}
European, paid at $T$ on $\phi(S_T - K) > 0$: cash-or-nothing (amount $A$) or asset-or-nothing ($S_T$).
\emph{Black--Scholes}, vol at the strike: $D^{\mathrm{cash}} = A\,P(0,T)\,N(\phi d_2)$, $D^{\mathrm{asset}} = P(0,T)\,F\,N(\phi d_1)$.
This ignores the smile slope. \emph{On the smile} the unit cash digital is $-\phi\,\partial V/\partial K$ of the vanilla
priced at its own vol, $P(0,T)\,N(\phi d_2) - \phi\,\mathcal{V}\,\partial\sigma/\partial K$ ($\mathcal{V}$ the vanilla
vega), computed as the centred spread $\phi\,[V(K - hK) - V(K + hK)]/(2hK)$; asset-or-nothing $= K\,D^{\mathrm{cash}}_{A=1}
+ \phi\,V(K)$. With a negative skew the smile lowers the digital put and raises the digital call; the gap is first
order (one vega times the skew).
\emph{Monte Carlo}: the indicator has no pathwise derivative, so it is replaced by the sigmoid
$1/(1 + e^{-\phi(S_T - K)/(\varepsilon K)})$ of width $\varepsilon$ (1\% of the strike by default; $\varepsilon = 0$ is the exact indicator, with
no pathwise delta or vega).

% ─────────────────────────────────────────────────────────────────────────────
\subsection{Barrier options}\label{sec:mp-barrier}
European call or put, knocked in or out when the spot reaches the barrier $H$ from above (down) or from below (up);
the barrier is monitored continuously or on given dates (a single date at expiry is the barrier observed at expiry
only); no rebate. Knock-out $=$ vanilla $-$ knock-in.

\paragraph{Black--Scholes} (one vol: the surface vol at $(K, T)$):
\begin{itemize}[nosep]
  \item continuous monitoring: the closed form of \citet{reiner1991} with carry $b$ \citep{haug2007};
  \item at expiry only: exact by vanillas and cash digitals struck at $H$, e.g.\ the down-and-in put is
        $\mathrm{Put}(H) + (K - H)\,\mathrm{DigitalPut}(H)$ for $K > H$;
  \item $n$ dates: the continuous formula at the shifted barrier $H\,e^{\pm 0.5826\,\sigma\sqrt{T/n}}$ of \citet{broadie1997} (moved away
        from the spot).
\end{itemize}

\paragraph{Monte Carlo.} Continuous monitoring by the Brownian-bridge probability of not crossing the barrier between
two simulated points \citep{beaglehole1997,glasserman2003}, multiplied over the steps:
\[
  \prod_i \Big(1 - e^{-2\,\ln(S_i/H)^+ \ln(S_{i+1}/H)^+/(\sigma_H^2\,\Delta t_i)}\Big)\quad(\text{down; } \ln(H/S) \text{ for up}),
\]
$\sigma_H$ the surface vol at $(H, T)$: exact for a lognormal model with that vol, and continuous in the path (pathwise
sensitivities). Discrete monitoring checks the barrier on the dates, the indicator smoothed as for the digitals.

% ─────────────────────────────────────────────────────────────────────────────
\subsection{Quanto payment}\label{sec:mp-quanto}
A quanto product pays its payoff, computed in the underlying's currency $f$, in another currency $d$ at a fixed rate
(one unit of payoff = one unit of $d$): a EUR-paid option on SPX, a USD-paid autocallable on SX5E. With $X$ the value of
one unit of $f$ in $d$, the underlying's dynamics under the measure of $d$ has an extra drift
\[
  \frac{dS}{S} = \big(r_f - q - \rho_{S,X}\,\sigma_S\,\sigma_X\big)\,dt + \sigma_S\,dW,
\]
and the payments are discounted on the curve of $d$. $\rho_{S,X}$ is the correlation of the underlying with $X$ (the
stored pair, or minus its correlation when the stored pair is the inverse, e.g.\ EURUSD for a USD underlying paid in
EUR) and $\sigma_X$ the vol of $X$; both come from the market data (correlation matrix, FX vol). Black--Scholes: the
forward becomes $F\,e^{-\rho\,\sigma_S\,\sigma_X T}$, with $\sigma_S$ the vol at the strike and the discount factor of
$d$ (vanilla, digital, barrier). Local vol: the drift uses the local vol of the path,
$-\rho\,\sigma_{\mathrm{loc}}(t, S_t)\,\sigma_X\,dt$, so the quanto forward has no closed form and the paths are not
rescaled to it. The FX vol is, for now, the historical vol of the pair (52 weekly returns, annualised) and the
correlation the median of its 13-week rolling correlations with the underlying; the greeks reach both by AAD.

\emph{Composite.} The payoff on the underlying converted into $d$, $Y = S\,X$, with strikes and levels in $d$ (a
EUR-paid put on SPX struck in EUR). $Y$ is the price in $d$ of a traded asset, so its forward is $F_S\,F_X(T)$ (with
$F_X$ the forward of $X$ by covered interest parity) and, with lognormal $S$ and $X$, its vol is
$\sigma_Y^2 = \sigma_S^2 + \sigma_X^2 + 2\rho\,\sigma_S\sigma_X$: Black--Scholes on $Y$, discounted in $d$, $\sigma_S$ at
the strike in the underlying's currency $K / F_X(T)$. Monte Carlo: $S$ with the quanto drift above (its dynamics
under the measure of $d$) and $X_t = F_X(t)\,e^{-\sigma_X^2 t/2 + \sigma_X W^X_t}$ on the same draws,
$W^X = \rho\,W^S + \sqrt{1-\rho^2}\,Z$; the payoff on $S_t X_t$. A payoff in the underlying's currency converted at
the future FX rate needs no model: its value is the price in that currency converted at the FX spot (the price
currency of the pricers).

\subsection{Autocallable}\label{sec:mp-autocallable}
A strike $K$ (the barriers in \% of it); coupon fixings $t_i$ with rates $c_i$, barriers $B^c_i$ and payment dates $p_i$, with or
without memory; autocall fixings $u_j$ with barriers $B^a_j$ (a step-down lowers them by a step $\delta$ at each
observation down to a floor $B^a_{\min}$: $B^a_j = \max(B^a - j\,\delta, B^a_{\min})$, $j = 0, 1, \dots$); at the last
coupon fixing a put knocked in below $B^p$ (observed
at that date only), with strike $K_p$ (by default $K$) and participation $\pi$. Valued on the investor's side, coupons minus put, per unit notional; on each
path, with $\alpha$ = 1 until the first autocall and $m$ the coupons unpaid so far:
\[
  V = N\,\E\Big[\sum_i \alpha_{t_i}\,\mathbb 1\{S_{t_i} \ge B^c_i\}\,c_i(1 + m_{t_i})\,P(0,p_i)
  \;-\; \alpha_T\,\mathbb 1\{S_T < B^p\}\,\pi\,\frac{(K_p - S_T)^+}{K}\,P(0,p_n)\Big],
\]
$m$ reset when a coupon is paid and increased by one when it is missed ($m \equiv 0$ without memory); $\alpha$ set to 0 when
$S_{u_j} \ge B^a_j$. \emph{Forward start}: with a strike date $t_0$ after today, $K = S_{t_0}$ is fixed on each path and every
barrier and $K_p$ are in \% of it; the payoff depends on $S_t / S_{t_0}$ only, so the price hardly depends on today's spot
(no delta with a lognormal model and no cash dividends) and is driven by the smile after $t_0$ (the forward smile).
Monte Carlo on a model of \S\ref{sec:mp-models} (local vol by default), on the fixing dates (and $t_0$) merged with a
grid of daily steps over the first three months and weekly ones after: a local vol builds the smile of a short
maturity only through steps small against it (weekly steps leave the 1--3 month wings of an index several vol
points low). Each indicator is a sigmoid of width $\varepsilon K$ so that delta and vega exist pathwise.

\emph{Live trades.} After the trade date some observations are already fixed: their values are the official closes
of the underlyings on those dates (or the fixings of the trade's confirmation). The value at $t$ is the conditional
expectation given them: the strike is the close of the strike date, and the past observations enter the formula
above as numbers instead of simulated values, with exact indicators (they are facts, not smoothed). They set the
state of the paths: coupons paid or missed (the memory $m$), autocall already triggered ($\alpha = 0$: only what is
fixed and not yet paid is left). Payments before $t$ are settled and drop out; a coupon fixed but paid after $t$ is
certain and discounted. A fixing on $t$ itself is the spot (today's close is not known yet). The past fixings do not
move with the market: the greeks come from the future observations only.

\emph{Vega by expiry and skew.} The vega per node of a local vol price alternates in sign along an expiry: one node
moves the slope and the curvature of the smile, so the local vol is raised at its strike and lowered next to it
(a fine grid of nodes makes this worse, \citealp[ch.~13]{savine2018}). The surface enters the price only through
its nodes, so the first-order change for a deformation $\Delta\sigma_{ij}$ of the nodes is
$\sum_{ij} \mathcal V_{ij}\,\Delta\sigma_{ij}$, with no repricing. Two shapes per expiry $T_i$ are reported:
the level, $\Delta\sigma_{ij} = 1$ (all its strikes up one vol point), and the skew, the smile rotated around the
forward, $\Delta\sigma_{ij} = (1 - m_{ij})/10\%$ with $m_{ij} = K_j / F(T_i)$ held within $50\%$--$150\%$ (puts up,
calls down: a steeper skew). On a 3-year phoenix (coupon barrier 70\%, put below 60\%) the level and the skew of the
whole surface match central bumps on the same paths within $0.2\%$ and $2\%$ for a one-point move; the skew is convex
(the second-order part is $10\%$ of the first at one point, $20\%$ at two). The skew of a single expiry is exact to
first order but not linear already at a fraction of a point: rotating one row only bends the surface in time, the
calendar derivative of the local vol changes sign in the wings and its floor acts.

\emph{Worst-of.} On underlyings $S^1, \dots, S^n$ with initial levels $K_1, \dots, K_n$ every level is in \% of the worst
performance $P_t = \min_k S^k_t / K_k$: the formula above with $S_t / K$ replaced by $P_t$ (coupon and autocall
indicators, knock-in and put $\pi (K_p/K - P_T)^+$). The underlyings follow their own models (local vol or GBM, each on
its market forward) driven by Brownian motions with the correlation matrix of the market, $dW^k\,dW^l = \rho_{kl}\,dt$
(independent normals times the Cholesky factor of $\rho$); an underlying in another currency than the product's has
the quanto drift of \S\ref{sec:mp-quanto}. The payoff falls when any underlying falls, so the price increases with the
correlations (an investor sold a worst-of is long correlation); the derivative of the minimum goes to the worst
underlying of each path, so the AAD gives a delta and a vega per node for each underlying and a sensitivity to each
correlation pair. With the autocall out
of reach and no put, the product is a strip of digitals; with no coupons, a put knocked in at expiry
(\S\ref{sec:mp-digital}, \S\ref{sec:mp-barrier}).

\emph{Side.} The value is the client's (receives the coupons, sells the put) or the issuer's (pays the coupons, buys
the put): the same number with the other sign, the cash flows too.

\emph{Autocallable swap.} The autocallable against a funding leg: the client also pays an index plus a margin $m$ on
periods $[s_i, e_i]$ from the strike date to each coupon payment date, until the autocall. The rate is not simulated
(rates and equity independent): period $i$ pays $(F_i + m)\,\tau_i$ at $e_i$, $F_i = (P_I(s_i)/P_I(e_i) - 1)/\tau_i$
the forward of the index curve $P_I$ and $\tau_i$ the ACT/360 year fraction, discounted on the payment currency's
curve and weighted by the probability that the product is alive at $s_i$:
\[ V_{\text{funding}} = N \sum_i (F_i + m)\,\tau_i\,P(0, e_i)\,\Q(\text{not called by } s_i), \]
an autocall paid on a date ends the funding there (the period ending on it is paid). The probabilities come from the
same paths and smoothed indicators as the coupons, so delta, vega and skew include the termination of the funding;
without an autocall the leg is the floating leg of a swap. A period already started takes the forward from today
(no rate fixings in the market yet).
